%% ============================================================
%%  SECCIÓN 7 — Resultados Experimentales
%% ============================================================

\section{Resultados Experimentales}
\label{sec:resultados}

En esta sección se presentan los parámetros instrumentales fijos de operación y el cuadro completo de datos crudos registrados para las 9 combinaciones posibles entre los 3 resistores y los 3 condensadores del módulo experimental en la Mesa C7.

En la Tabla~\ref{tab:parametros_constantes} se detallan las magnitudes que permanecieron constantes durante toda la adquisición de datos en el osciloscopio digital.

\begin{table}[H]
    \centering
    \caption{Parámetros instrumentales constantes de la excitación periódica.}
    \label{tab:parametros_constantes}
    \resizebox{\linewidth}{!}{%
    \begin{tabular}{lccc}
        \toprule
        \textbf{Magnitud física} & \textbf{Símbolo} & \textbf{Valor nominal} & \textbf{Incertidumbre asociada} \\
        \midrule
        Frecuencia de la onda cuadrada & $f$ & \qty{250,0}{\hertz} & $\pm\qty{1}{\ppm}$ (\qty{\pm 0,00025}{\hertz}) \\
        Período de la onda & $T = 1/f$ & \qty{4,000}{\milli\second} & $\pm\qty{0,05}{\micro\second}$ \\
        Semiperíodo de carga / descarga & $T/2$ & \qty{2,000}{\milli\second} & $\pm\qty{0,05}{\micro\second}$ \\
        Tensión pico a pico del generador & $V_0$ & \qty{10,00}{\voltpp} & $\pm\qty{0,05}{\volt}$ \\
        Ciclo de trabajo (\texttt{Duty Cycle}) & $D$ & \qty{50,0}{\percent} & $\pm\qty{0,1}{\percent}$ \\
        Resistencia interna del generador & $R_{\text{gen}}$ & \qty{50}{\ohm} & $\pm\qty{1}{\ohm}$ \\
        \bottomrule
    \end{tabular}%
    }
\end{table}

En la Tabla~\ref{tab:datos_experimentales} se presentan las 9 configuraciones ensayadas sistemáticamente en nuestro puesto de trabajo (Mesa C7). Cada combinación fue implementada en el módulo de ensayo conectando sucesivamente cada uno de los tres resistores con los tres condensadores disponibles, registrando en cada caso la respuesta transitoria en la pantalla del osciloscopio digital.

Para cada ensayo se reporta la resistencia $R$ con su incertidumbre estándar combinada $u(R)$ ($k=1$), la frecuencia $f = \qty{250}{\hertz}$, la constante de tiempo experimental $\tau_{\text{exp}}$ medida con cursores digitales al \qty{63,2}{\percent} de saturación junto a su incertidumbre estándar $u(\tau)$ ($k=1$), la capacitancia experimental calculada $C_{\text{exp}} = \tau / R$ con su incertidumbre expandida $U(C) = 2 \cdot u_c(C)$ ($k=2$, nivel de confianza del \qty{95,45}{\percent}), la capacitancia nominal rotulada $C_{\text{nom}}$ y el error relativo porcentual respecto a dicha referencia de fábrica.

\begin{table}[H]
    \centering
    \caption{Matriz experimental completa de constantes de tiempo y capacitancias para la Mesa C7.}
    \label{tab:datos_experimentales}
    \resizebox{\linewidth}{!}{%
    \begin{tabular}{ccccccc}
        \toprule
        \textbf{Fila} & $R \pm u(R)\ (k=1)$ & $f$ & $\tau_{\text{exp}} \pm u(\tau)\ (k=1)$ & $C_{\text{exp}} \pm U(C)\ (k=2)$ & $C_{\text{nom}}$ & \textbf{Error} \\
         & (\unit{\kilo\ohm}) & (\unit{\hertz}) & (\unit{\micro\second}) & (\unit{\nano\farad}) & (\unit{\nano\farad}) & (\unit{\percent}) \\
        \midrule
        \textbf{1} & \num{9,980 \pm 0,058} & \num{250} & \num{107,8 \pm 0,14} & \num{10,80 \pm 0,13} & \num{10,00} & $+8{,}02$ \\
        \textbf{2} & \num{9,980 \pm 0,058} & \num{250} & \num{273,6 \pm 0,14} & \num{27,41 \pm 0,32} & \num{22,40} & $+22{,}39$ \\
        \textbf{3} & \num{9,980 \pm 0,058} & \num{250} & \num{344,8 \pm 0,15} & \num{34,55 \pm 0,40} & \num{33,23} & $+3{,}97$ \\
        \textbf{4} & \num{6,832 \pm 0,033} & \num{250} & \num{74,2 \pm 0,14} & \num{10,86 \pm 0,11} & \num{10,00} & $+8{,}61$ \\
        \textbf{5} & \num{6,832 \pm 0,033} & \num{250} & \num{187,0 \pm 0,14} & \num{27,37 \pm 0,27} & \num{22,40} & $+22{,}19$ \\
        \textbf{6} & \num{6,832 \pm 0,033} & \num{250} & \num{237,9 \pm 0,14} & \num{34,82 \pm 0,34} & \num{33,23} & $+4{,}79$ \\
        \textbf{7} & \num{2,250 \pm 0,012} & \num{250} & \num{24,4 \pm 0,14} & \num{10,84 \pm 0,17} & \num{10,00} & $+8{,}44$ \\
        \textbf{8} & \num{2,250 \pm 0,012} & \num{250} & \num{61,5 \pm 0,14} & \num{27,33 \pm 0,31} & \num{22,40} & $+22{,}02$ \\
        \textbf{9} & \num{2,250 \pm 0,012} & \num{250} & \num{77,5 \pm 0,14} & \num{34,44 \pm 0,38} & \num{33,23} & $+3{,}65$ \\
        \bottomrule
    \end{tabular}%
    }
\end{table}

En la Tabla~\ref{tab:datos_experimentales} se observa con claridad que para un mismo condensador (por ejemplo, el condensador $C_1$ en las filas 1, 4 y 7), la constante de tiempo $\tau$ se reduce de manera estrictamente monótona conforme disminuye el valor resistivo: desde $\tau_1 = \qty{107,8}{\micro\second}$ con $R_1 \approx \qty{10}{\kilo\ohm}$, pasando por $\tau_4 = \qty{74,2}{\micro\second}$ con $R_2 \approx \qty{6,8}{\kilo\ohm}$, hasta $\tau_7 = \qty{24,4}{\micro\second}$ con $R_3 \approx \qty{2,25}{\kilo\ohm}$. No obstante, el cociente experimental $C = \tau / R$ se mantiene notablemente estable en torno a \qty{10,8}{\nano\farad}, evidenciando la constancia física de la capacitancia intrínseca independientemente del resistor acoplado a la malla.